\section{Lorentz Group Digression}

The Lorentz group is defined as the set of all $4$x$4$ matrices which preserve the scalar products 
with respect to the relativistic metric, that is, the set of all $\Lambda:x^\mu\rightarrow {x'}^\mu=\tensor{\Lambda}{^\mu_\nu}x^\nu$, 
such,
\begin{align}
    {x}^\mu g_{\mu\nu} {y}^\nu&={x'}^\alpha g_{\alpha\beta} {y'}^\beta\\
    {x}^\mu g_{\mu\nu} {y}^\nu&=\tensor{\Lambda}{^\alpha_\mu}{x}^\mu g_{\alpha\beta} {y}^\nu\tensor{\Lambda}{^\beta_\nu}\\
    {x}^\mu g_{\mu\nu} {y}^\nu&={x}^\mu {y}^\nu \tensor{\Lambda}{^\alpha_\mu}g_{\alpha\beta} \tensor{\Lambda}{^\beta_\nu}
\end{align}
as this should be valid for all $x,y$,
\begin{align}
     g_{\mu\nu} &= \tensor{{\Lambda^{\textnormal{T}}}}{_\mu^\alpha}g_{\alpha\beta} \tensor{\Lambda}{^\beta_\nu}\\
     g&=\Lambda^{\textnormal{T}}g\Lambda
\end{align}
in other words, the Lorentz group is the group of all orthogonal $4$x$4$ matrices with respect to the non-positive definite 
metric $g$. For orthogonal $n$x$n$ matrices with respect to the standard $\mathbb R^n$ metric, this group is denoted $O\qty(n)$, 
as we're in a $1$ negative and $3$ positive eigenvalues metric, we denote by $O\qty(1,3)\cong O\qty(3,1)$.
Apart from the topological structure carried by $O\qty(1,3)$, we are mainly interested in the tangent space structure at 
the origin, at the identity element, also know as the Lie algebra, $\mathfrak o\qty(1,3)$. This is the case because 
the generators of the Lie algebra are in a one to one correspondence with hermitian operators in the Hilbert space of 
the quantum theory which represents conserved charges, for the Lorentz group these are the angular momentum and boost 
operators. It's not hard to obtain the connected piece of $O\qty(1,3)$, we have two conditions which splits our group 
into $4$ connected pieces, 
\begin{align}
    \Det\qty[g]&=\Det[\Lambda^{\textnormal{T}}]\Det[g]\Det[\Lambda]\\
    -1&=-\Det[\Lambda^{\textnormal{T}}]\Det[\Lambda]\\
    1&=\qty(\Det[\Lambda])\Rightarrow \Det[\Lambda] =\pm 1
\end{align}
and also,
\begin{align}
    g_{00}&=\tensor{{\Lambda^{\textnormal{T}}}}{_0^\alpha}g_{\alpha\beta}\tensor{\Lambda}{^\beta_0}\\
    -1&=\tensor{{\Lambda^{\textnormal{T}}}}{_0^0}g_{00}\tensor{\Lambda}{^0_0}+\tensor{{\Lambda^{\textnormal{T}}}}{_0^i}g_{ij}\tensor{\Lambda}{^j_0}\\
    \tensor{\Lambda}{^0_0}\tensor{\Lambda}{^0_0}&=1+\sum\limits_{i}\norm{\tensor{\Lambda}{^i_0}}^2\geq 1\\
    \qty(\tensor{\Lambda}{^0_0})^2&\geq 1\Rightarrow \tensor{\Lambda}{^0_0}\begin{cases}\geq 1\\\leq -1\end{cases}
\end{align}

The connected piece is the one satisfying $\Det[\Lambda]= +1$ and $\tensor{\Lambda}{^0_0}\geq 1$, which is called the 
proper orthochronous Lorentz group $SO^+\qty(1,3)$, proper ($S$) for $\Det[\Lambda]=+1$ and orthochronous ($\.^+$) for 
$\tensor{\Lambda}{^0_0}\geq 1$. It's not hard to see that from the connected piece we can get all the other three pieces by 
compositions with the following matrices,
\begin{align}
    \tensor{\mathcal P}{^\mu_\nu}=\textnormal{diag}\mqty(1&-1&-1&-1)&,\ \ \tensor{\mathcal T}{^\mu_\nu}=\textnormal{diag}\mqty(-1&1&1&1)\\
    \tensor{\qty[\mathcal P\mathcal T]}{^\mu_\nu}&=\textnormal{diag}\mqty(-1&-1&-1&-1)
\end{align}
in fact they form a group by themselves, $\qty{\mathbbm 1,\mathcal P,\mathcal T,\mathcal P\mathcal T}$, called the Klein four group. With them we 
can reconstruct the full Lorentz group, 
\begin{align}
    O\qty(1,3) = \qty{\mathbbm 1,\mathcal P,\mathcal T,\mathcal P\mathcal T}\rtimes SO^+\qty(1,3)
\end{align}

The quantum theory allows for projective representations of groups, this is due to states of the Hilbert space which differ by a phase represent 
the same physical configuration, hence, representations of groups can obey the group multiplication law up to a phase, a.k.a. projective representations. 
Constructing all the projective representations is similar to constructing all the proper representations of a cover of the group, for quantum mechanics 
the double cover is enough. It's known that the double cover of $SO^+(1,3)$ is $SL(2,\mathbb C)$, the group of $2$x$2$ matrices of unit determinant over the 
complex numbers, that is, \begin{align}SO^+\qty(1,3)\cong SL(2,\mathbb C)/\mathbb Z_2\end{align} the real quantum symmetry group is then $SL(2,\mathbb C)$, due 
to it being the double cover, it contains proper spinorial representations, which are not present at $SO^+(1,3)$, at least as proper representations.

Let's construct this isomorphism, first, a disclaimer, we'll use a little bit of a posteriori knowledge, which could be obtained through the item $1.b$ without 
knowing of this construction. What we need is a way to represent Lorentz vectors over the $2$x$2$ complex values matrices, in which a proper orthochronous Lorentz transformation 
is given through a unit determinant transformation. The correct choice, which we'll leave to explain why latter,
\begin{align}
    X = \mqty(-x^0+x^3& x^1-\im x^2\\x^1+\im x^2& -x^0-x^3)=-x^0\mathbbm 1 +x^i\sigma_i = x^\mu\sigma_\mu,\ \ \ \Det[X] = -x_\mu x^\mu
\end{align}
with $\sigma_i$ being the usual Pauli matrices, we also defined $\sigma_ 0 =-\mathbbm 1$, we leave for now open as if this is an honest four vector, 
but in fact it transforms covariantly over $SL(2,\mathbb C)$, and even more, it's invariant. As the scalar product is given by the determinant, 
Lorentz transformations are linear matrix multiplications which preserve the determinant,
\begin{align}
    X' = \mqty(-{x'}^0+{x'}^3& {x'}^1-\im {x'}^2\\{x'}^1+\im {x'}^2& -{x'}^0-{x'}^3)= \tilde U X \tilde V,\ \ \ \Det[X']=\Det[X]
\end{align}
also, the transformation needs to preserve the reality of each of the components, this can be assured by noting that $X^\dagger  =X$, and then 
demanding that ${X'}^\dagger = X'$, which fixes $\tilde V ={\tilde U}^\dagger$. The preservation of the determinant gives the following constraint,
\begin{align}
    \Det[X]&=\Det[X']\\
    \Det[X]&=\Det[\tilde U]\Det[X]\Det[{\tilde U}^\dagger]\\
    1&=\norm{\Det[\tilde U]}^2\Rightarrow \Det[\tilde U]=\exp(\im \alpha)
\end{align}
notice that $\tilde U$ and $U=\exp(-\im\frac\alpha2)\tilde U$ generates the same transformation on $X$, hence, we can consider only the determinant one 
matrices. What we concluded up to now is that at least the $2$x$2$ unit determinant matrices over the complexes is surjective in the (proper orthochronous?) Lorentz 
group, let's verify it's the proper orthochronous, to see the condition $\tensor{\Lambda}{^0_0}\geq 1$, take $X = -\mathbbm 1\leftrightarrow x=\textnormal{diag}\mqty(1&0&0&0)$, then, 
\begin{align}
    X'=\mqty(-\tensor{\Lambda}{^0_0}+\tensor{\Lambda}{^3_0}& \tensor{\Lambda}{^1_0}-\im \tensor{\Lambda}{^2_0}\\\tensor{\Lambda}{^1_0}+\im \tensor{\Lambda}{^2_0}& -\tensor{\Lambda}{^0_0}-\tensor{\Lambda}{^3_0})\Rightarrow \Tr[X'] = -2\tensor{\Lambda}{^0_0}
\end{align}
thus,
\begin{align}
    \tensor{\Lambda}{^0_0}&=-\frac12\Tr[X']=-\frac12\Tr[UXU^\dagger]=\frac12\Tr[UU^\dagger]
\end{align}
as $\Det[U]=1$, we have $\Det[UU^\dagger]=1$, if $\lambda_1,\lambda_2$ are the eigenvalues of $U$, $\lambda_1^\ast,\lambda_2^\ast$ are the eigenvalues of 
$U^\dagger$, hence, the eigenvalues of $UU^\dagger$ are $\norm{\lambda_1}^2,\norm{\lambda_2}^2>0$, the constraint $\Det[UU^\dagger]= 1 =\norm{\lambda_1}^2+\norm{\lambda_2}^2$ 
imposes that $\Tr[UU^\dagger] = \norm{\lambda_1}^2+\norm{\lambda_1}^2\geq 2$, therefore,
\begin{align}
    \tensor{\Lambda}{^0_0}&=\frac12\Tr[UU^\dagger]\geq 1
\end{align}
this clears that $SL(2,\mathbb C)$ is at least surjective in the orthochronous Lorentz group. Now, the proper part, it suffices to show that there isn't in 
this group a parity transformation. Suppose exists $U_p$ such that,
\begin{align}
    X' &= U_pXU_p'^\dagger=U_p\qty( -x^0 \mathbbm 1+x^i\sigma_i)U_p^\dagger = -x^0 \mathbbm 1-x^i\sigma_i
\end{align}
which imposes,
\begin{align}
    U_pU^\dagger_p&=\mathbbm 1\\
    U_p\sigma_iU^\dagger_p&=-\sigma_i
\end{align}

Now, $\sigma_1\sigma_2\sigma_3 = \im\mathbbm 1$ (check it!), thus, we can compute the following,
\begin{align}
    U_p\sigma_1U^\dagger_pU_p\sigma_2U^\dagger_pU_p\sigma_3U^\dagger_p&=(-)^3\sigma_1\sigma_2\sigma_3\\
    U_p\sigma_1\sigma_2\sigma_3U^\dagger_p&=-\im\mathbbm 1\\
    U_p\im \mathbbm 1U^\dagger_p&=-\im\mathbbm 1\\
    U_pU^\dagger_p&=-\mathbbm 1
\end{align}
contradiction! Thus the parity transformation with $\Det[\mathcal P]=-1$ isn't present in $SL(2,\mathbb C)$, and as this is a connected continuous group, 
none of the transformations with $\Det[\Lambda]=-1$ are present. This concludes that $SL(2,\mathbb C)$ is surjective over the proper orthochronous Lorentz 
group. But is still up to discussion whether this is an injective map, the answer is, it isn't, more than one transformation in $SL(2,\mathbb C)$ correspond 
to the same $SO^+(1,3)$ transformation. This can be seen from the transformation law,
\begin{align}
    X' = UXU^\dagger= (-U)X(-U)^\dagger
\end{align}
as $U$ is two dimensional, $\Det[-U]=1$, and we conclude that both $U$ and $-U$ generate the same proper orthochronous Lorentz transformation, formally this means 
that $SO^+(1,3)$ can be seen as an equivalence class over $SL(2,\mathbb C)$, 
\begin{align}
    SO^+ (1,3) \cong \qty{U\in SL(2,\mathbb C)| U\sim -U}
\end{align}
the action $U\rightarrow -U$ can be seen as a $\mathbb Z_2$ action over $SL(2,\mathbb C)$, hence, this equivalence class is a quotient group,
\begin{align}
    SO^+ (1,3) \cong  SL(2,\mathbb C)/\mathbb Z_ 2
\end{align}
this also means that $\mathfrak{o} \qty(1,3)\cong \mathfrak{so} \qty(1,3)\cong\mathfrak{so}^+ \qty(1,3) \cong \mathfrak {sl}(2,\mathbb C)$. The presence 
of $\mathbb Z_ 2$ instead of $\mathbb Z_n$ signals that $SL(2,\mathbb C)$ is the double-cover, actually in this case it is the universal-cover. For general 
dimensions the double-cover of $SO^+(p,q)$ is defined as $\textnormal{Spin}^+(p,q)$,
\begin{align}
    SO^+ (p,q) \cong  \textnormal{Spin}^+(p,q)/\mathbb Z_ 2
\end{align}
in a comically similar fashion, we have the following,
\begin{align}
    SO (p,q) \cong  \textnormal{Spin}(p,q)/\mathbb Z_ 2,\ \ \ O (p,q) \cong  \textnormal{Pin}(p,q)/\mathbb Z_ 2
\end{align}
these groups $\textnormal{Spin}^+(p,q),\textnormal{Spin}(p,q),\textnormal{Pin}(p,q)$ are more easily definable by means of Clifford algebras, 
which we'll see soon enough, a representation of the Clifford algebra $\textnormal{Cl}(p,q)$ is given by the Dirac gamma matrices. The interesting point is 
that $\textnormal{Spin}(1,3)\cong \textnormal{Spin}(3,1)$ contains the $\mathcal P\mathcal T$ transformation, but not parity and time inversion separately, 
but the Pin group, which contains all $\mathcal P,\mathcal T,\mathcal P\mathcal T$ has the incredible property of
\begin{align}
    \textnormal{Pin}(1,3)\ncong\textnormal{Pin}(3,1)
\end{align}
which is a consequence of a non isomorphism between the Clifford algebras,
\begin{align}
    \textnormal{Cl}(1,3)\ncong\textnormal{Cl}(3,1)
\end{align}

This is simply outstanding, which in principle means that in a quantum theory with fermions and in which all $\mathcal P,\mathcal T,\mathcal P\mathcal T$ are 
true symmetries there could be noticeable experimental results that could decide in favor of $\textnormal{Pin}(3,1)$ over $\textnormal{Pin}(1,3)$.